\Chapter{CONCLUSION}\label{sec:Conclusion}
Texte / Text.

%%
%%  SYNTHESE DES TRAVAUX / SUMMARY OF WORKS
%%
\section{Synthèse des travaux / Summary of Works}
Texte / Text.

%%
%%  LIMITATIONS
%%
\section{Limitations de la solution proposée / Limitations}\label{sec:Limitations}

Les conditions sous lesquelles le certificat du filtre de sécurité est établi, ainsi que les quantités surveillées à l'exécution qui en signalent la violation, sont énoncées à la Section~\ref{subsec:garanties}. Cette section ne les répète pas : elle rassemble ce que le système, une fois ces conditions tenues, ne sait pas faire.

\paragraph{Sécurité sans progression.} La garantie obtenue porte sur la sécurité, non sur l'accomplissement de la tâche. L'invariance de l'ensemble sûr vaut pour une géométrie d'obstacle arbitraire, tandis que le dégagement postural (Section~\ref{subsec:minima}) n'étend la vivacité qu'aux blocages portés par le corps du bras, résolubles dans le noyau de la tâche. Lorsque l'obstacle barre le chemin du lien contrôlé lui-même, cas de la paroi frontale et de la cavité concave faisant face à la trajectoire nominale, la projection annule la composante entrante de la vitesse, le modèle génératif réémet la même intention à chaque replanification, et le système demeure indéfiniment en arrêt sûr : la barrière reste positive, mais la tâche ne progresse plus. Cette limitation est inhérente à un filtre qui opère par projection locale sans planification globale. La résolution des impasses topologiques relève de la couche de planification nominale, dont le remplacement par une politique consciente des obstacles, ou l'adjonction d'un routage global, sort du cadre de ce travail. Les situations de blocage ainsi induites sont caractérisées au Chapitre~\ref{sec:Theme3}.

\paragraph{Marge dimensionnée pour la simulation.} La marge de sécurité $d_{\mathrm{safe}}$ retenue est propre au cadre d'évaluation : les nuages d'obstacles y sont échantillonnés sur la surface exacte des objets et les transformations de repères sont connues, de sorte que les termes d'erreur perceptive, bruit de profondeur, calibration et quantification de carte, y sont nuls. Un déploiement sur matériel exigerait d'augmenter cette marge du budget d'erreur perceptive correspondant, et de revalider les constantes qui en dépendent.

\paragraph{Invariance certifiée sur une barrière apprise.} Le certificat porte sur la barrière apprise, non sur la géométrie réelle : la distance physique hérite de l'erreur d'approximation du champ de distance. En régime établi, la marge conserve un dégagement réel confortable devant cette erreur, mais dans le pire transitoire certifié le plancher résiduel ne domine plus l'erreur de queue du modèle, et la conjonction des deux pires cas ne peut être exclue par le seul calcul. L'évitement physique est donc validé empiriquement, par la distance mesurée au maillage exact, et non déduit du certificat.

%%
%%  AMELIORATIONS FUTURES / FUTURE RESEARCH
%%
\section{Améliorations futures / Future Research}

\paragraph{Du bouclier géométrique à la gestion du contact.} Le filtre proposé traite la sécurité comme un problème d'évitement : il maintient une marge strictement positive entre le robot et tout ce qu'il perçoit. La tâche d'alimentation assistée exige pourtant, dans sa phase terminale, un contact volontaire avec l'utilisateur, que cette formulation interdit par construction (Section~\ref{subsec:revue_cbf_qp}). Lever cette incompatibilité suppose d'adjoindre au bouclier géométrique un second niveau de sécurité, fondé sur l'estimation en ligne des efforts externes et sur une transition vers une commande en impédance ou en admittance lors de l'insertion, la marge $d_{\mathrm{safe}}$ devenant alors fonction de l'état de la tâche plutôt que constante. Ce niveau relève de l'interaction physique humain-robot et des seuils physiologiques de la norme ISO/TS 15066:2016, exclus de la portée de ce travail, et constitue la suite naturelle du système sur une plateforme réelle.
